\begin{center}
\textbf{\large{Tóm tắt}}
\end{center}

\addcontentsline{toc}{chapter}{Tóm tắt}

\begin{small}

Báo cáo trình bày quá trình phân tích, thiết kế và xây dựng ứng dụng di động \textbf{Fitty} nhằm hỗ trợ người dùng thiết lập kế hoạch tập luyện và theo dõi sức khỏe cá nhân hóa trên nền tảng Android. Bài toán đặt ra là xây dựng một hệ thống không chỉ dừng ở việc hiển thị danh sách bài tập, mà còn phải hỗ trợ toàn bộ luồng sử dụng từ xác thực người dùng, onboarding thu thập hồ sơ ban đầu, sinh kế hoạch tập khởi đầu, tra cứu nội dung bài tập, thực hiện buổi tập, theo dõi tiến độ và duy trì tương tác qua thông báo.

Để giải quyết bài toán, hệ thống được xây dựng theo kiến trúc nhiều tầng với Kotlin và Jetpack Compose ở tầng giao diện, use case và repository ở tầng nghiệp vụ, Room và DataStore ở tầng lưu trữ cục bộ, cùng Firebase làm hạ tầng phía sau cho xác thực, lưu trữ dữ liệu đám mây, lưu trữ media và thông báo đẩy. Một điểm kỹ thuật nổi bật của hệ thống là việc áp dụng mô hình \textit{offline-first} cho phân hệ bài tập, trong đó metadata bài tập được đồng bộ từ Firestore, chuẩn hóa và lưu vào Room để giao diện có thể truy cập nhanh và ổn định hơn. Ngoài ra, hệ thống còn hỗ trợ quản lý nội dung động thông qua Firebase Content, cho phép cập nhật onboarding, danh mục bài tập, mẫu kế hoạch khởi đầu và cấu hình hành vi mà không cần sửa mã giao diện.

Kết quả đạt được của báo cáo là một nguyên mẫu chức năng tương đối đầy đủ cho ứng dụng hỗ trợ tập luyện cá nhân hóa. Hệ thống đã hiện thực được các phân hệ chính gồm: đăng ký và đăng nhập, dùng thử ở chế độ khách, onboarding và xem trước kế hoạch, quản lý kế hoạch tập, tra cứu bài tập theo nhóm cơ, xem chi tiết và media bài tập, thực hiện buổi tập nhanh hoặc buổi tập theo lịch, theo dõi tiến độ cơ bản, quản lý hồ sơ người dùng, thông báo cục bộ và thông báo đẩy FCM. Bên cạnh đó, dự án cũng đã có một số kiểm thử đơn vị cho các thành phần nghiệp vụ quan trọng như sinh kế hoạch khởi đầu, chính sách đồng bộ bài tập, hoàn thành buổi tập và theo dõi tiến độ.

Nhìn chung, báo cáo cho thấy việc kết hợp kiến trúc Android hiện đại với các dịch vụ đám mây và lưu trữ cục bộ là một hướng phù hợp để phát triển các ứng dụng sức khỏe có tính cá nhân hóa. Kết quả đạt được là nền tảng tốt để tiếp tục mở rộng hệ thống theo hướng nâng cao chất lượng gợi ý kế hoạch, tăng cường trực quan hóa tiến độ và hoàn thiện hơn trải nghiệm người dùng trong thực tế.

\vspace*{1cm}
\textbf{Từ khóa}: Android, Jetpack Compose, Firebase, offline-first, ứng dụng tập luyện cá nhân hóa

\end{small}
